\section{Governing Equation and Spatial Discretization}
\label{sec:governing_equation}

\subsection{Dimensionless Energy Equation}

The objective of this project is to develop a numerical solution
for the two-dimensional energy equation of an incompressible flow.
The equation describes the evolution of temperature under the
combined effects of advection, diffusion, and viscous dissipation.
Its dimensionless form is given by
\begin{equation}
\begin{aligned}
\frac{\partial T}{\partial t}
+\frac{\partial (uT)}{\partial x}
+\frac{\partial (vT)}{\partial y}
={}&
\frac{1}{\mathrm{Re}\,\mathrm{Pr}}\nabla^2 T
\\
&+
\frac{\mathrm{Ec}}{\mathrm{Re}}
\left[
2\left(\frac{\partial u}{\partial x}\right)^2
+2\left(\frac{\partial v}{\partial y}\right)^2
+\left(
\frac{\partial u}{\partial y}
+\frac{\partial v}{\partial x}
\right)^2
\right],
\end{aligned}
\label{eq:energy}
\end{equation}
where $T$ is the dimensionless temperature, $u$ and $v$ are the
velocity components in the $x$- and $y$-directions, respectively.
The dimensionless parameters $\mathrm{Re}$, $\mathrm{Pr}$, and
$\mathrm{Ec}$ denote the Reynolds, Prandtl, and Eckert numbers.

The left-hand side of Equation~\eqref{eq:energy} contains the
temporal and advective contributions. The first term on the
right-hand side represents thermal diffusion, while the second
term accounts for viscous dissipation, which represents the
conversion of mechanical energy into thermal energy due to
velocity gradients. This contribution is treated as a source
term, denoted by $S$.

\subsection{Finite-Volume Formulation}

The governing equation is discretized using the finite-volume
method (FVM). The computational domain is divided into
rectangular control volumes, with the solution variables stored
at their centers. Integrating Equation~\eqref{eq:energy} over
a control volume and applying the divergence theorem to the
advective and diffusive terms converts the corresponding volume
integrals into surface integrals over the control-volume faces.

For a two-dimensional rectangular control volume, the surface
integrals are evaluated as contributions from the east, west,
north, and south faces. The resulting semi-discrete equation is
\begin{equation}
\begin{aligned}
\frac{d\overline{T}_{ij}}{dt}
={}&
-\frac{F^a_e-F^a_w}{\Delta x}
-\frac{F^a_n-F^a_s}{\Delta y}
\\
&+
\frac{1}{\mathrm{Re}\,\mathrm{Pr}}
\left[
\frac{F^d_e-F^d_w}{\Delta x}
+
\frac{F^d_n-F^d_s}{\Delta y}
\right]
+\overline{S}_{ij},
\end{aligned}
\label{eq:fv_semidiscrete}
\end{equation}
where $\overline{T}_{ij}$ and $\overline{S}_{ij}$ are the
control-volume averages of temperature and the source term.
The quantities $F^a$ and $F^d$ denote the normal advective
flux and the normal temperature gradient, respectively.
The subscripts $e$, $w$, $n$, and $s$ represent the east,
west, north, and south faces. The control-volume dimensions
are $\Delta x$ and $\Delta y$.

The surface integrals are approximated using values at the
face centers. Since the numerical solution is stored at
cell centers, the face values are reconstructed from the
available cell-centered quantities. Second-order central
approximations are employed for the spatial discretization.

\subsection{Spatial Discretization of the Advective Terms}

The energy equation contains both advective and diffusive
terms, whose different physical characteristics must be
considered when selecting the spatial discretization scheme.
Diffusion spreads information in all directions and is
therefore discretized using a second-order central scheme.
Although central differencing can also be used for advection,
it may produce nonphysical oscillations or instability in
convection-dominated regions, particularly when the cell
Péclet number is large. In such regions, an upwind scheme
is generally more appropriate because it respects the direction
of information transport.

In the present work, second-order central differencing is used
for the advective terms. This choice is appropriate provided
that diffusion is sufficiently strong relative to convection.

The advective terms are written in conservative form as
spatial derivatives of the products $uT$ and $vT$. In the
present discretization, each product is treated as a single
quantity. For a generic cell-centered variable $\phi$, the
first-order spatial derivatives are approximated using
second-order central differences:
\begin{align}
\left.\frac{\partial \phi}{\partial x}\right|_{ij}
&\approx
\frac{\phi_{i+1,j}-\phi_{i-1,j}}{2\Delta x},
\label{eq:first_derivative_x}
\\
\left.\frac{\partial \phi}{\partial y}\right|_{ij}
&\approx
\frac{\phi_{i,j+1}-\phi_{i,j-1}}{2\Delta y}.
\label{eq:first_derivative_y}
\end{align}

Applying these expressions directly to the products $uT$
and $vT$ gives
\begin{equation}
\left.\frac{\partial (uT)}{\partial x}\right|_{ij}
\approx
\frac{
u_{i+1,j}T_{i+1,j}
-u_{i-1,j}T_{i-1,j}
}{2\Delta x},
\label{eq:advection_x}
\end{equation}
and
\begin{equation}
\left.\frac{\partial (vT)}{\partial y}\right|_{ij}
\approx
\frac{
v_{i,j+1}T_{i,j+1}
-v_{i,j-1}T_{i,j-1}
}{2\Delta y}.
\label{eq:advection_y}
\end{equation}

For completeness, the cell-centered temperature can also be
interpolated to the control-volume faces using second-order
central interpolation:
\begin{align}
T_{i+\frac{1}{2},j}
&\approx \frac{T_{i,j}+T_{i+1,j}}{2},
&
T_{i-\frac{1}{2},j}
&\approx \frac{T_{i-1,j}+T_{i,j}}{2},
\label{eq:temperature_interpolation_x}
\\
T_{i,j+\frac{1}{2}}
&\approx \frac{T_{i,j}+T_{i,j+1}}{2},
&
T_{i,j-\frac{1}{2}}
&\approx \frac{T_{i,j-1}+T_{i,j}}{2}.
\label{eq:temperature_interpolation_y}
\end{align}
The same central interpolation can be applied to other
cell-centered variables when face values are required.
However, Equations~\eqref{eq:advection_x} and
\eqref{eq:advection_y} define the specific advective
discretization adopted here: central differences are applied
directly to the products $uT$ and $vT$.

\subsection{Spatial Discretization of the Diffusive Term}

The diffusive contribution is discretized using second-order
central differences. The first derivatives of temperature are
approximated as
\begin{align}
\left.\frac{\partial T}{\partial x}\right|_{ij}
&\approx
\frac{T_{i+1,j}-T_{i-1,j}}{2\Delta x},
\\
\left.\frac{\partial T}{\partial y}\right|_{ij}
&\approx
\frac{T_{i,j+1}-T_{i,j-1}}{2\Delta y}.
\end{align}

The second derivatives are approximated by
\begin{align}
\left.\frac{\partial^2 T}{\partial x^2}\right|_{ij}
&\approx
\frac{T_{i+1,j}-2T_{i,j}+T_{i-1,j}}{\Delta x^2},
\label{eq:second_derivative_x}
\\
\left.\frac{\partial^2 T}{\partial y^2}\right|_{ij}
&\approx
\frac{T_{i,j+1}-2T_{i,j}+T_{i,j-1}}{\Delta y^2}.
\label{eq:second_derivative_y}
\end{align}

Consequently, the Laplacian of temperature is discretized as
\begin{equation}
\begin{aligned}
\nabla^2 T\big|_{ij}
\approx{}&
\frac{T_{i+1,j}-2T_{i,j}+T_{i-1,j}}{\Delta x^2}
\\
&+
\frac{T_{i,j+1}-2T_{i,j}+T_{i,j-1}}{\Delta y^2}.
\end{aligned}
\label{eq:discrete_laplacian}
\end{equation}

\subsection{Spatial Discretization of the Viscous-Dissipation Source}

The viscous-dissipation source term is evaluated using
velocity gradients at the cell centers. The derivatives are
approximated using second-order central differences:
\begin{align}
\left.\frac{\partial u}{\partial x}\right|_{ij}
&\approx
\frac{u_{i+1,j}-u_{i-1,j}}{2\Delta x},
\\
\left.\frac{\partial v}{\partial y}\right|_{ij}
&\approx
\frac{v_{i,j+1}-v_{i,j-1}}{2\Delta y},
\\
\left.\frac{\partial u}{\partial y}\right|_{ij}
&\approx
\frac{u_{i,j+1}-u_{i,j-1}}{2\Delta y},
\\
\left.\frac{\partial v}{\partial x}\right|_{ij}
&\approx
\frac{v_{i+1,j}-v_{i-1,j}}{2\Delta x}.
\end{align}

Substituting these expressions into the source-term equation
gives
\begin{equation}
\begin{aligned}
S_{ij}
=
\frac{\mathrm{Ec}}{\mathrm{Re}}
\Bigg[
&2\left(
\frac{u_{i+1,j}-u_{i-1,j}}{2\Delta x}
\right)^2
\\
&+2\left(
\frac{v_{i,j+1}-v_{i,j-1}}{2\Delta y}
\right)^2
\\
&+\left(
\frac{u_{i,j+1}-u_{i,j-1}}{2\Delta y}
+
\frac{v_{i+1,j}-v_{i-1,j}}{2\Delta x}
\right)^2
\Bigg].
\end{aligned}
\label{eq:discrete_viscous_source}
\end{equation}

\subsection{Complete Spatially Discretized Equation}

Combining the central-difference approximations for the
advective and diffusive terms with the discretized
viscous-dissipation source gives the complete
spatially discretized energy equation:
\begin{equation}
\boxed{
\begin{aligned}
\frac{d\overline{T}_{ij}}{dt}
={}&
-\frac{1}{2\Delta x}
\left(
u_{i+1,j}T_{i+1,j}
-u_{i-1,j}T_{i-1,j}
\right)
\\[0.15cm]
&+
\frac{1}{\mathrm{Re}\,\mathrm{Pr}}
\frac{
T_{i+1,j}-2T_{i,j}+T_{i-1,j}
}{\Delta x^2}
\\[0.15cm]
&-
\frac{1}{2\Delta y}
\left(
v_{i,j+1}T_{i,j+1}
-v_{i,j-1}T_{i,j-1}
\right)
\\[0.15cm]
&+
\frac{1}{\mathrm{Re}\,\mathrm{Pr}}
\frac{
T_{i,j+1}-2T_{i,j}+T_{i,j-1}
}{\Delta y^2}
+S_{ij}
\\[0.15cm]
={}& FI_{ij}+S_{ij}
=RHS_{ij}.
\end{aligned}
}
\label{eq:complete_spatial_discretization}
\end{equation}

Here, $S_{ij}$ denotes the discretized viscous-dissipation
source term. Equation~\eqref{eq:complete_spatial_discretization}
is the spatially discretized, semi-discrete energy equation.
All spatial derivatives have been replaced by algebraic
expressions, whereas the temporal derivative remains
continuous. The resulting system of ordinary differential
equations can then be advanced in time using an appropriate
time-integration method. The temporal discretization and
the numerical solution methods used in this project are
discussed in the next section.

For the parameter study, the spatially discretized energy
equation is reformulated in terms of the Peclet number $Pe$,
the Brinkman number $Br$, and the convection coefficient
$c$. Using the definitions
\[
Pe=\mathrm{Re}\,\mathrm{Pr},
\qquad
Br=Ec\,\mathrm{Pr},
\]
the coefficient of the viscous-dissipation term becomes
\[
\frac{Ec}{\mathrm{Re}}=\frac{Br}{Pe}.
\]
The parameterized equation therefore takes the form:
\begin{equation}
\boxed{
\begin{aligned}
\frac{d\overline{T}_{ij}}{dt}
={}&
-\frac{c}{2\Delta x}
\left(
u_{i+1,j}T_{i+1,j}
-u_{i-1,j}T_{i-1,j}
\right)
\\[0.15cm]
&-\frac{c}{2\Delta y}
\left(
v_{i,j+1}T_{i,j+1}
-v_{i,j-1}T_{i,j-1}
\right)
\\[0.15cm]
&+\frac{1}{Pe}
\left[
\frac{T_{i+1,j}-2T_{i,j}+T_{i-1,j}}{\Delta x^2}
\right.
\\
&\qquad\left.
+\frac{T_{i,j+1}-2T_{i,j}+T_{i,j-1}}{\Delta y^2}
\right]
+\frac{Br}{Pe}\,\Phi_{ij}
\\[0.15cm]
={}&FI_{ij}(Pe,c)
+\frac{Br}{Pe}\,\Phi_{ij}
=RHS_{ij},
\end{aligned}
}
\label{eq:parameterized_energy}
\end{equation}
where $\Phi_{ij}$ denotes the discretized
viscous-dissipation function without its multiplicative
coefficient. The operator $FI_{ij}(Pe,c)$ contains the
advective and diffusive contributions. This formulation
provides the basis for the parameter study, in which the
effects of convection, thermal diffusion, and viscous
dissipation are examined through controlled variations of
$c$, $Pe$, and $Br$.

\subsection{Implementation of Boundary Conditions Using Ghost Cells}
\label{subsec:ghost_cell_boundary_conditions}

The spatial discretization requires temperature values at neighboring
cells, including near the boundaries of the computational domain.
However, the physical domain does not contain cells outside its
boundaries. To apply the same discrete spatial operators to boundary
cells, an additional layer of fictitious cells, known as
\textit{ghost cells}, is introduced around the physical domain.

The temperature values in ghost cells are not independent unknowns.
Instead, they are calculated from the prescribed boundary conditions
and the temperature values in adjacent physical cells. This approach
allows the discretized equations to retain a consistent stencil near
the boundaries.

In the following expressions, the physical cells are indexed by
$i=1,\ldots,n_x$ and $j=1,\ldots,n_y$. The inlet, bottom-wall,
top-wall, and outlet ghost cells are located at
$i=0$, $j=0$, $j=n_y+1$, and $i=n_x+1$, respectively.

\subsubsection{Dirichlet Boundary Conditions}

For a prescribed temperature at a boundary face, the ghost-cell
temperature is determined by assuming that the boundary-face
temperature is the average of the adjacent physical-cell and
ghost-cell temperatures. At the inlet, this relation is
\begin{equation}
T_{\mathrm{in},j}
=
\frac{T_{0,j}+T_{1,j}}{2},
\end{equation}
which gives
\begin{equation}
\boxed{
T_{0,j}=2T_{\mathrm{in},j}-T_{1,j}.
}
\label{eq:ghost_inlet}
\end{equation}

Similarly, the bottom and top walls have prescribed temperatures
$T_{\mathrm{bottom}}$ and $T_{\mathrm{top}}$. The corresponding
ghost-cell relations are
\begin{equation}
\boxed{
T_{i,0}
=
2T_{\mathrm{bottom}}-T_{i,1},
}
\label{eq:ghost_bottom}
\end{equation}
and
\begin{equation}
\boxed{
T_{i,n_y+1}
=
2T_{\mathrm{top}}-T_{i,n_y}.
}
\label{eq:ghost_top}
\end{equation}

For the present problem, the inlet temperature varies linearly
with the wall-normal coordinate, while the bottom and top wall
temperatures are zero and one, respectively:
\begin{equation}
T_{\mathrm{in}}(y)=y,
\qquad
T_{\mathrm{bottom}}=0,
\qquad
T_{\mathrm{top}}=1.
\end{equation}
Consequently, the ghost-cell values simplify to
\begin{align}
T_{0,j}&=2y_j-T_{1,j},
\label{eq:ghost_inlet_specific}
\\
T_{i,0}&=-T_{i,1},
\label{eq:ghost_bottom_specific}
\\
T_{i,n_y+1}&=2-T_{i,n_y}.
\label{eq:ghost_top_specific}
\end{align}

\subsubsection{Outlet Boundary Condition}

At the outlet, a zero streamwise temperature gradient is imposed:
\begin{equation}
\left.\frac{\partial T}{\partial x}\right|_{x=L_x}=0.
\label{eq:outlet_zero_gradient}
\end{equation}
Using a first-order approximation to the gradient between the
last physical cell and its ghost cell gives
\begin{equation}
\frac{T_{n_x+1,j}-T_{n_x,j}}{\Delta x}=0.
\end{equation}
Therefore, the outlet ghost-cell temperature is
\begin{equation}
\boxed{
T_{n_x+1,j}=T_{n_x,j}.
}
\label{eq:ghost_outlet}
\end{equation}

The ghost-cell relations are applied whenever the spatial operator
is evaluated. They provide the temperature values required by the
discrete equations near the boundaries without introducing
additional independent unknowns. In the subsequent numerical
solution methods, the same boundary conditions are retained, while
their algebraic contributions are incorporated according to the
structure of each solver.

\subsubsection{Boundary Conditions in Delta Form}

In the approximate factorization method, the unknown is the
temperature increment between consecutive time levels, defined as
\begin{equation}
\delta T_{i,j}^{n+1}
=
T_{i,j}^{n+1}-T_{i,j}^{n}.
\label{eq:delta_temperature}
\end{equation}
The ghost-cell relations must therefore also be expressed in terms
of these increments. Since the prescribed inlet and wall
temperatures are constant in time, their increments at the
boundary faces are zero. Applying the ghost-cell relations at
successive time levels and subtracting them gives the corresponding
incremental boundary conditions.

At the inlet,
\begin{equation}
\begin{aligned}
T_{0,j}^{n+1}&=2T_{\mathrm{in},j}-T_{1,j}^{n+1},\\
T_{0,j}^{n}&=2T_{\mathrm{in},j}-T_{1,j}^{n},
\end{aligned}
\end{equation}
and consequently,
\begin{equation}
\boxed{
\delta T_{0,j}^{n+1}
=
-\delta T_{1,j}^{n+1}.
}
\label{eq:delta_ghost_inlet}
\end{equation}

Similarly, for the bottom and top walls,
\begin{equation}
\boxed{
\delta T_{i,0}^{n+1}
=
-\delta T_{i,1}^{n+1},
}
\label{eq:delta_ghost_bottom}
\end{equation}
and
\begin{equation}
\boxed{
\delta T_{i,n_y+1}^{n+1}
=
-\delta T_{i,n_y}^{n+1}.
}
\label{eq:delta_ghost_top}
\end{equation}

At the outlet, the zero-gradient condition gives
\begin{equation}
T_{n_x+1,j}^{n+1}=T_{n_x,j}^{n+1},
\qquad
T_{n_x+1,j}^{n}=T_{n_x,j}^{n},
\end{equation}
so that
\begin{equation}
\boxed{
\delta T_{n_x+1,j}^{n+1}
=
\delta T_{n_x,j}^{n+1}.
}
\label{eq:delta_ghost_outlet}
\end{equation}

These incremental relations allow the spatial operator to be
evaluated consistently in terms of the unknown temperature updates.
They also ensure that the prescribed inlet and wall temperatures
remain unchanged after each update, while the zero-gradient
condition continues to hold at the outlet.